\begin{itemframe}{ضرب موازی ماتریس‌ها}
\itm
در این بخش، بررسی می‌کنیم که چگونه می‌توان الگوریتم‌های ضرب ماتریس را موازی‌سازی کرد.
خواهیم دید که هر الگوریتم را می‌توان به‌سادگی با استفاده از کلید‌واژه‌هایی که آموختیم موازی‌سازی نمود.
\itm
این الگوریتم‌ها را با استفاده از کار-گستره تحلیل خواهیم کرد و مشاهده می‌کنیم که الگوریتم‌های موازی همان کارایی الگوریتم سری متناظر خود را روی یک پردازنده به دست می‌آورند، در حالی که روی تعداد زیادی پردازنده نیز مقیاس‌پذیر هستند.
\end{itemframe}


\begin{itemframe-s}{ضرب موازی ماتریس‌ها}{موازی‌سازی روش حلقه‌های تودرتو}
\itm
اولین الگوریتمی که بررسی می‌کنیم
\texttt{P-MATRIX-MULTIPLY}
است که به‌سادگی دو حلقهٔ بیرونی در ضرب ماتریسی عادی را موازی‌سازی می‌کند.
\begin{algorithm}[H]\alglr
\caption{P-MATRIX-MULTIPLY}
\begin{algorithmic}[1]
\Func{P-MATRIX-MULTIPLY}{$A, B, C, n$}
\pFor{$i = 1 \ \textbf{to} \ n$} \Comment{compute entries in each of $n$ rows}
    \pFor{$j = 1 \ \textbf{to} \ n$} \Comment{compute $n$ entries in row $i$}
        \For{$k = 1 \ \textbf{to} \ n$}
            \State $c_{ij} \gets c_{ij} + a_{ik} \cdot b_{kj}$
        \EndFor
    \EndpFor
\EndpFor
\end{algorithmic}
\end{algorithm}

\end{itemframe-s}


\begin{itemframe-s}{ضرب موازی ماتریس‌ها}{موازی‌سازی روش حلقه‌های تودرتو}
\itm
اکنون این الگوریتم را تحلیل می‌کنیم. از تحلیل تصویر سری این الگوریتم در می‌یابیم کار برابر با است با:
$$
T_1(n) = \Theta(n^3)
$$

\itm
زیرا مسیر بحرانی شامل حرکت بازگشتی ناشی از حلقهٔ موازی خط 1، سپس درخت بازگشتی ناشی از حلقهٔ موازی خط 2، و سپس اجرای تمام n تکرار حلقهٔ معمولی خط 3 است. در نتیجه گستره برابر می‌شود با:
$$
T_\infty(n) = \Theta(\lg n) + \Theta(\lg n) + \Theta(n) = \Theta(n)
$$
بنابراین، نرخ موازی‌سازی برابر است با:
$$
\Theta(n^3) / \Theta(n) = \Theta(n^2).
$$
\end{itemframe-s}


\begin{itemframe-s}{ضرب موازی ماتریس‌ها}{موازی‌سازی روش تقسیم و حل}
\itm
احتمالاً با انجام ضرب ماتریسی با استفاده از تقسیم و حل آشنایی داشته باشید. در این قسمت می‌خواهیم همان روش را موازی‌سازی کنیم.
\itm
این روش دو ماتریس
$n \times n$
را در زمان
$\Theta(n^3)$
ضرب می‌کند.
\itm
بیایید ببینیم چگونه می‌توان این الگوریتم را با استفاده از فراخوانی‌ها موازی بازگشتی پیاده سازی کرد.
\end{itemframe-s}


\begin{itemframe-s}{ضرب موازی ماتریس‌ها}{موازی‌سازی روش تقسیم و حل}
\itm
می‌دانیم که تابع سریالی ضرب ماتریس به روش تقسیم و حل، سه ماتریس
$n \times n$
دریافت می‌کند
و سپس محاسبهٔ ماتریسی زیر را انجام می‌دهد: (فرض کنید سه ماترسی را A B C بنامیم)
$$
C = C + A \cdot B
$$
این کار با انجام هشت ضرب از زیر ماتریس‌های
$\frac{n}{2} \times \frac{n}{2}$
از A و B به‌صورت بازگشتی انجام می‌شود.
\itm

تابع
\texttt{P-MATRIX-MULTIPLY-RECURSIVE}
در اسلاید بعدی همین راهبرد تقسیم‌ و حل را پیاده‌سازی می‌کند، اما برای انجام این هشت ضرب از فراخوانی موازی استفاده می‌کند.
برای اجتناب از بروز عدم قطعیت رقابتی هنگام به‌روزرسانی عناصر، یک ماتریس موقتی D ساخته می‌شود و سپس C و D جمع می‌شوند تا نتیجهٔ نهایی به دست آید.
\end{itemframe-s}


\begin{itemframe-s}{ضرب موازی ماتریس‌ها}{موازی‌سازی روش تقسیم و حل}
\itm
الگوریتم ضرب ماتریسی موازی به روش تقسیم و حل به این صورت است:
 \begin{algorithm}[H]\alglr
\caption{P-MATRIX-MULTIPLY-RECURSIVE}
\begin{algorithmic}[1]
\Func{P-MATRIX-MULTIPLY-RECURSIVE}{$A, B, C, n$}
\If{$n = 1$} \Comment{base case}
    \State $c_{11} \gets c_{11} + a_{11} \cdot b_{11}$
    \State \Return
\EndIf
\State let $D$ be a new $n \times n$ matrix \Comment{temporary matrix}
\pFor{ $i = 1$ to $n$}
    \pFor{$j = 1$ to $n$}
        \State $d_{ij} \gets 0$
    \EndFor
\EndFor
\State partition $A, B, C, D$ into $\tfrac{n}{2} \times \tfrac{n}{2}$ submatrices:
\State \hspace{1.5em} $A_{11}, A_{12}, A_{21}, A_{22}$,
$B_{11}, B_{12}, B_{21}, B_{22}$,

\State \hspace{1.5em} $C_{11}, C_{12}, C_{21}, C_{22}$,
and $D_{11}, D_{12}, D_{21}, D_{22}$
\end{algorithmic}
\end{algorithm}

\end{itemframe-s}


\begin{itemframe-s}{ضرب موازی ماتریس‌ها}{موازی‌سازی روش تقسیم و حل}
\itm
ادامه تابع ضرب ماتریسی موازی به روش تقسیم و حل...
 \begin{algorithm}[H]\alglr
\begin{algorithmic}[1]
\setcounter{ALG@line}{10}
\State \textbf{spawn} P-MATRIX-MULTIPLY-RECURSIVE($A_{11}, B_{11}, C_{11}, n/2$)
\State \textbf{spawn} P-MATRIX-MULTIPLY-RECURSIVE($A_{11}, B_{12}, C_{12}, n/2$)
\State \textbf{spawn} P-MATRIX-MULTIPLY-RECURSIVE($A_{21}, B_{11}, C_{21}, n/2$)
\State \textbf{spawn} P-MATRIX-MULTIPLY-RECURSIVE($A_{21}, B_{12}, C_{22}, n/2$)
\State \textbf{spawn} P-MATRIX-MULTIPLY-RECURSIVE($A_{12}, B_{21}, D_{11}, n/2$)
\State \textbf{spawn} P-MATRIX-MULTIPLY-RECURSIVE($A_{12}, B_{22}, D_{12}, n/2$)
\State \textbf{spawn} P-MATRIX-MULTIPLY-RECURSIVE($A_{22}, B_{21}, D_{21}, n/2$)
\State \textbf{spawn} P-MATRIX-MULTIPLY-RECURSIVE($A_{22}, B_{22}, D_{22}, n/2$)
\State \textbf{sync} \Comment{wait for spawned submatrix products}
\For{parallel $i = 1$ to $n$} \Comment{update $C = C + D$}
    \For{parallel $j = 1$ to $n$}
        \State $c_{ij} \gets c_{ij} + d_{ij}$
    \EndFor
\EndFor
\end{algorithmic}
\end{algorithm}
\end{itemframe-s}


\begin{itemframe-s}{ضرب موازی ماتریس‌ها}{موازی‌سازی روش تقسیم و حل}
\itm
خطوط 1 تا 3 حالت پایه را که مربوط به ضرب ماتریس‌های
$1 \times 1$
است، مدیریت می‌کنند و بخش باقی‌مانده با حالت بازگشتی سروکار دارد. خط 8 هر یک از چهار ماتریس
$A, B, C, D$
را به زیرماتریس‌های
$
\frac{n}{2} \times \frac{n}{2}
$
تقسیم می‌کند.
\itm
در خطوط ۱۱ تا ۱۸ هشت ضرب ماترسی مذکور به شکل موازی انجام شده و در نتیجه در ماتریس‌ها C‌ و D قرار می‌گیرد.
\itm
در خط ۱۹ دستور sync تضمین می‌کند که همهٔ ضرب‌ها محاسبه شده‌باشند، و پس از آن در خطوط ۲۰ تا ۲۲ ماتریس C و D جمع می‌شوند.
\end{itemframe-s}


\begin{itemframe-s}{ضرب موازی ماتریس‌ها}{موازی‌سازی روش تقسیم و حل}
\itm
بیایید رویهٔ
\texttt{P-MATRIX-MULTIPLY-RECURSIVE}
را تحلیل کنیم. ابتدا، کار یا همان
$M_1(n)$
را تحلیل می‌کنیم.
\itm
می‌دانیم باید قسمتی که هشت ضرب بازگشتی
را انجام می‌دهد از بقیهٔ قسمت‌ها پیچیدگی زمانی بالاتری دارد و پیچیدگی زمانی ‌آن برابر است با:
$$
M_1(n) = 8M_1(n/2) + \Theta(n^2) = \Theta(n^3)
$$

که با حالت 1 قضیهٔ اصلی
\fn{master theorem}
به دست می‌آید.
\end{itemframe-s}


\begin{itemframe-s}{ضرب موازی ماتریس‌ها}{موازی‌سازی روش تقسیم و حل}
\itm
اکنون باید طول مسیر بحرانی یا همان
$M_\infty(n)$
را تعیین کنیم. از آنجا که هشت فراخوانی بازگشتی موازی همگی روی ماتریس‌هایی با اندازهٔ یکسان اجرا می‌شوند، گستره همه فراخوانی‌های بازگشتی با هم برابر است و بنابراین بیشینه گستره آنها برابر است با گستره یکی از آنها و برابر است با:
$$
M_\infty(n/2)
$$
\itm
گستره برای حلقه‌های تو در توی موازی در خطوط 5 تا 7 برابر با
$\Theta(\lg n)$
است، به‌طور مشابه، حلقه‌های تو در توی موازی در خطوط ۲۰ تا ۲۲ نیز
$\Theta(\lg n)$
دیگر به گستره اضافه می‌کنند.
پیچیدگی زمانی تقسیم ماتریس هم ثابت است و از آن صرف نظر می‌کنیم.
 بنابراین گستره این تابع به صورت زیر است:

$$M_\infty(n) = M_\infty(n/2) + \Theta(\lg n)$$
\end{itemframe-s}


\begin{itemframe-s}{ضرب موازی ماتریس‌ها}{موازی‌سازی روش تقسیم و حل}
\itm
این رابطه بازگشتی تحت حالت دوم قضیهٔ اصلی با قرار می‌گیرد و جواب آن چنین است:
$$
M_\infty(n) = \Theta(\lg^2 n)
$$
\itm
پس نرخ موازی‌سازی الگوریتم
\texttt{P-MATRIX-MULTIPLY-RECURSIVE}
برابر است با:
$$
M_1(n) / M_\infty(n) = \Theta(n^3 / \lg^2 n)
$$

که بسیار بزرگ است و نشان دهنده تأثیر بسزای موازی سازی در افزایش سرعت این الگوریتم است.
\end{itemframe-s}